% Paper skeleton based on the course's IEEE template (MVSR Paper Template WS2022).
% Rules: main part (first page up to and including references) must fill EXACTLY three pages.
% Figures: serif font, >= body text size, >= 300 DPI or vector PDF, readable in grayscale.
% Every claim must be backed by literature or by our results.
\documentclass[conference]{IEEEtran}
\usepackage{cite}
\usepackage{amsmath,amssymb,amsfonts}
\usepackage{graphicx}
\usepackage{textcomp}
\usepackage{xcolor}
\usepackage{subcaption}
\usepackage{booktabs}
\usepackage{url}

\newcommand{\todo}[1]{\textcolor{red}{[TODO: #1]}}

\begin{document}

\title{Comparing SLAM Systems Across Sensor Modalities\\on a Shared Ground-Robot Dataset}

\author{\IEEEauthorblockN{Elias Bitsch}
\and
\IEEEauthorblockN{Philip Stix}
\and
\IEEEauthorblockN{Viktoriia Ovdiienko}
\and
\IEEEauthorblockN{Jiayi Zhou}
}

\maketitle

\begin{abstract}
% No citations, symbols, maths or footnotes. Must state: motivation, contribution, verifiability.
\todo{Motivation (indoor service robots, changing light, people). Contribution: reproducible comparison of four SLAM systems with different sensor modalities on identical M3DGR sequences, using a standard metric (ATE). Main finding in one sentence. Code and results are publicly available (repository link).}
\end{abstract}

\section{Introduction}
% ~0.3 pages. Why SLAM on ground robots, why compare sensor modalities, contribution in 2-3 points.
\todo{Motivation, problem, contribution.}

\section{State of the Art}
% ~0.5 pages. Terminology, existing solutions, "so what" -> our gap.
SLAM is commonly formulated as maximum-a-posteriori estimation over a factor graph~\cite{cadena2016past,grisetti2010tutorial}.
\todo{One paragraph per evaluated system (each member writes their own):}
\todo{Elias: MASt3R-SLAM~\cite{murai2025mast3rslam}.}
\todo{Philip: <system>.}
\todo{Viktoriia: <system>.}
\todo{Jiayi: <system>.}
\todo{Existing comparisons, e.g.~\cite{trejos2022slam,zhang2025m3dgr}; what is missing.}

\section{Materials and Methods}
% ~0.7 pages. Dataset, systems + configuration, metrics, hardware, pipeline (cite software).
\subsection{Dataset}
\todo{M3DGR~\cite{zhang2025m3dgr}: robot, sensors, motion-capture ground truth, the three sequences and why.}

\subsection{Evaluated Systems}
\todo{Table: system, sensors, paradigm, loop closure, alignment, hardware. Default parameters only.}

\subsection{Evaluation Metrics}
We report the absolute trajectory error (ATE)~\cite{sturm2012benchmark}, computed with evo~\cite{grupp2017evo}.
Trajectories are aligned with the Umeyama method~\cite{umeyama1991least}: SE(3) for metric systems and Sim(3) for monocular systems, whose scale is unobservable~\cite{zhang2018tutorial}.
\todo{N runs per system and sequence, median and IQR, success rate, tracked fraction~\cite{campos2021orbslam3}.}

\section{Experimental Results}
% ~0.8 pages. Required: trajectory overlay of ALL systems, >= one map per system.
\begin{figure}[!t]
    \centering
    % \includegraphics[width=\columnwidth]{figures/trajectory_overlay.pdf}
    \fbox{\parbox[c][4cm][c]{0.9\columnwidth}{\centering trajectory overlay (notebooks/02\_trajectories.ipynb)}}
    \caption{Estimated trajectories of all systems and ground truth on \todo{sequence}, top-down view after alignment.}
    \label{fig:overlay}
\end{figure}

\begin{table}[!t]
    \centering
    \caption{ATE RMSE [m], median (IQR) over N runs, and success rate per sequence.}
    \label{tab:results}
    % \input{tables/results.tex}
    \todo{generated by notebooks/02\_trajectories.ipynb}
\end{table}

\todo{Fig.~\ref{fig:overlay}, Tab.~\ref{tab:results}, maps, box plot.}

\section{Summary and Outlook}
% ~0.5 pages. Interpret failure modes in the context of theory, limitations, outlook.
\todo{Failure modes vs. theory (scale unobservability, darkness, geometric degeneracy, drift without loop closure). Limitations: software-synchronised sensors, default parameters, small N, different hardware per system. Outlook.}

\bibliographystyle{IEEEtran}
\bibliography{references}

\end{document}
